% -- Encoding UTF-8 without BOM
% -- XeLaTeX => PDF (BIBER)

\documentclass[]{cv-style}          % Add 'print' as an option into the square bracket to remove colours from this template for printing. 
                                    % Add 'espanol' as an option into the square bracket to change the date format of the Last Updated Text

\sethyphenation{french}{} % Add words between the {} to avoid them to be cut 

\begin{document}

\header{Camille }{SIMON}           % Your name
\lastupdated

%----------------------------------------------------------------------------------------
%	SIDEBAR SECTION  -- In the aside, each new line forces a line break
%----------------------------------------------------------------------------------------

\begin{aside}
%
\section{Coordonnées}
25 Rue de la Vire
14000 Caen
France
~
+33 6 69 98 44 07
camille.nelly.simon@gmail.com
~
12 septembre 1992
%
\section{Langues}
Français (langue maternelle)
Anglais C1 (fluent)
%
\section{Programmation}
Python, Java, C, C++
Javascript, Scala
PHP, CSS3 \& HTML5
MySQL, PostgreSQL, Oracle, MongoDB
\LaTeX{}
Git
%
\end{aside}

%----------------------------------------------------------------------------------------
%	SKILLS SECTION
%----------------------------------------------------------------------------------------

%\section{skills}
%  \vspace{-0.2cm}

%Skill 1, skill 2, skill 3, skill 4, skill 5.

%----------------------------------------------------------------------------------------
%	WORK EXPERIENCE SECTION
%----------------------------------------------------------------------------------------

\section{Expériences}

\begin{entrylist}
%------------------------------------------------
\entry
  {2022--2023}
  {Développeuse et cheffe de projet, Communication Internet and Networks Solutions (CINS), Caen}
  {CDD}
  {Développement de thèmes et modules pour Prestashop et Wordpress. Gestion d'un portefeuille clients et encadrement des besoins et évolutions techniques.}
%------------------------------------------------
\entry
  {2021--2022}
  {Enseignante, La Manu Campus du Havre}
  {CDD}
  {Enseignements technique niveau Bac+1, Bac+2 : HTML/CSS, PHP, Python, Algorithmique. Etude de faisabilité pour la création d'un Mastère IoT.}
%------------------------------------------------
\entry
  {2018--2021}
  {Ingénieur d'étude, Université Le Havre Normandie, Le Havre}
  {CDD}
  {Réalisation d'une preuve de concept pour HAROPA Port du Havre portant sur la sécurisation des chaînes logistiques distribuées à l'aide de \textit{Blockchain} et \textit{smart contracts}. Travaux dirigés par Cyrille Bertelle et co-encadrée par Claude Duvallet.}
%------------------------------------------------
\entry
  {2018--2021}
  {Université Le Havre Normandie et ISEL, Le Havre}
  {Mission d'enseignement}
  {Cours de PIX à destination des licence 1\up{ère} année de sciences et des 1\up{ère} année du cycle préparatoire, préparation au certificat de compétences numériques.}
%------------------------------------------------
\entry
  {2015--2021}
  {Cous particuliers}
  {Enseignement}
  {Cours de maths niveau 1\up{ère} ES, Terminale S (préparation au bac), Licence de maths 1\up{ère} année. Cours d’informatique niveau Licence de maths et Licence d’informatique 1\up{ère} et 2\up{ème} année.}
%------------------------------------------------
\entry
  {04--08 2018}
  {Onecub, Calais}
  {Stage de recherche}
  {Réalisation d’une étude de faisabilité de la blockchain dans le cadre de la gestion des données personnelles au regard du RGPD.}
%------------------------------------------------
\entry
  {2015--2017}
  {Tutorats en langages Java, C et C++}
  {Mission d'enseignement}
  {Cours de renforcement à destination des étudiants de 2\up{ème} et 3\up{ème} année de licence de maths et d’informatique à l’université du Havre.}
%------------------------------------------------
\entry
  {2014--2017}
  {Université Le Havre Normandie, Le Havre}
  {Mission de communication}
  {Présentation des formations en mathématiques et informatique de l’Université lors de salons étudiants.}

%-----------------------------------------------
\entry
  {04--06 2016}
  {HAROPA Port du Havre, Le Havre}
  {Stage de recherche}
  {Réalisation d’une bibliothèque ferroviaire au sein du logiciel de modélisation multi-agent AnyLogic.}
%-----------------------------------------------
\end{entrylist}

%----------------------------------------------------------------------------------------
%	EDUCATION SECTION
%----------------------------------------------------------------------------------------

\section{Formations}

\begin{entrylist}
\entry
  {2017--2018}
  {Master 2 IWOCS}
  {Université Le Havre Normandie}
  {Ingénierie Web, Objets Communicants et Systèmes complexes Parcours orientation recherche.}
%------------------------------------------------
\entry
  {2016--2017}
  {Master 1 MATIS}
  {Université Le Havre Normandie}
  {MAThématiques et Informatique des Systèmes complexes et distribués. Mention Assez bien}
%------------------------------------------------
\entry
  {2016}
  {Licence Informatique}
  {Université Le Havre Normandie}
  {Mention Assez bien}
%------------------------------------------------
\entry
  {2012}
  {DUT Information-Communication}
  {IUT du Havre}
  {Diplôme Universitaire Technologie. Parcours InfoNum, Information Numérique dans les Organisations}
%-----------------------------------------------
\end{entrylist}

%----------------------------------------------------------------------------------------
%	OTHER QUALIFICATIONS SECTION
%----------------------------------------------------------------------------------------

%\section{Publications}

%\vspace{-0.2cm}

%\textbf{Camille Simon}, Cyrille Bertelle, Claude Duvallet, Jérôme Besancenot. \italica{Traçabilité du suivi des matières dangereuses dans les smart ports grâce aux technologies Blockchain}. \textit{RIRL (Rencontres Internationales de la Recherche en Logistique et SCM), Octobre 2020, Le Havre, France.}

%----------------------------------------------------------------------------------------
%	AWARDS SECTION
%----------------------------------------------------------------------------------------

\section{Centres d'intérêts} 

\vspace{-0.2cm}

Animation d'ateliers pour la bibliothèque universitaire du Havre : Initiation à Latex (2019), Découverte des Linux et d'Ubuntu (2020) et Prise en main d'Inkscape (2021).\\
Réalisation des logos du LMAH (Laboratoire de Mathématiques Appliquées du Havre) et du Master IWOCS.\\
Joueuse de Go, jeu de plateau asiatique comparable aux échecs (rang : 12 kyu, niveau FFG : -1072).

%----------------------------------------------------------------------------------------

\end{document}